\chapter{Theoretische Grundlagen}
\label{chap:theoretische-grundlagen}

Die empirische Bewertung der Resilienz einer GitOps-Control-Plane erfordert zunächst eine theoretische Grundlage. Diese wird im Folgenden geschaffen, indem auf die Entwicklung von \gl{iac} hin zu deklarativen Zustandsmodellen, die formale Definition des GitOps-Betriebsmodells, die regelungstheoretische Funktionsweise der Reconciliation-Schleife sowie die Konzepte der Fehlertoleranz und des Chaos Engineerings eingegangen wird. Ein Blick auf die zugrundeliegenden Netzwerk- und Betriebssystemmechanismen verdeutlicht zudem die Kausalkette zwischen physischen Paketverlusten und dem Verhalten auf Anwendungsebene.

\section{Deklarative Infrastruktur und das GitOps-Paradigma}
\label{sec:deklarative-infrastruktur-gitops}

Um die Bereitstellung und den Betrieb komplexer Softwaresysteme reproduzierbar zu gestalten, hat sich das Konzept von \textit{Infrastructure as Code} (\gl{iac}) als grundlegendes Paradigma etabliert. Dabei werden sämtliche Infrastruktur- und Konfigurationsartefakte in maschinenlesbaren Definitionsdateien beschrieben, anstatt Systeme manuell oder interaktiv zu konfigurieren.\autocite{morris_iac_2020} Innerhalb von \gl{iac} sowie im Konfigurationsmanagement wird auf konzeptioneller Ebene strikt zwischen imperativen und deklarativen Ansätzen differenziert. Beide Ansätze verfolgen das Ziel einer automatisierten Systembereitstellung, wählen dafür aber unterschiedliche Steuerungsprinzipien.\autocite{burns_kubernetes_2022,morris_iac_2020}

\subsection{Imperative Steuerung versus deklarative Zustandsmodelle}
\label{subsec:imperativ-vs-deklarativ}

Der imperative Ansatz (prozedurales Modell) definiert eine explizite Abfolge von ausführbaren Einzelschritten, die ein System von seinem aktuellen Zustand in einen gewünschten Zielzustand überführen sollen. Klassische \gl{ci}-/\gl{cd}-Pipelines implementieren dieses Prinzip häufig über ein sogenanntes \textit{Push-Modell}: Nach einem erfolgreichen Build-Prozess führt ein externer Pipeline-Runner Kommandos oder Skripte (etwa \texttt{kubectl apply} oder \gl{ssh}-Befehle) direkt gegen die Zielumgebung aus.\autocite{beetz_gitops_2022,morris_iac_2020} Obwohl dieses Vorgehen bei einfachen Bereitstellungen intuitiv erscheint, weist es in verteilten Systemen strukturelle Schwächen auf. Da die Ausführung nur zum Zeitpunkt des Pipeline-Laufs stattfindet, bleiben nachträgliche manuelle Eingriffe oder schleichende Zustandsabweichungen im Zielsystem (\textit{Configuration Drift}) bis zum nächsten Deployment unbemerkt. Zudem erfordert das Push-Modell, dass weitreichende Schreib- und Administrationsberechtigungen für den Produktionscluster außerhalb der eigentlichen Laufzeitumgebung im \gl{ci}-System hinterlegt werden, was die Angriffsfläche der Infrastruktur vergrößert.\autocite{beetz_gitops_2022,opengitops_principles_2021}

Demgegenüber beschreibt das deklarative Modell ausschließlich den angestrebten Soll-Zustand (\textit{Desired State}) eines Systems in Form unveränderlicher Datenstrukturen (beispielsweise als \gl{yaml}-Manifeste), ohne die konkreten Ausführungsschritte zur Erreichung dieses Zustands vorzugeben.\autocite{burns_borg_2016,burns_kubernetes_2022} Die Verantwortung dafür, die Differenz zwischen dem aktuell beobachteten Ist-Zustand (\textit{Actual State} bzw. \textit{Live State}) und dem deklarierten Soll-Zustand zu berechnen und auszugleichen, wird vollständig an eine Laufzeitkomponente innerhalb des Zielsystems delegiert. Diese konsequente Trennung von Zustandsbeschreibung und Zustandsüberführung reduziert die Komplexität der Konfigurationsartefakte und bildet die Voraussetzung für selbstheilende Infrastrukturen.\autocite{burns_borg_2016,beyer_sre_2016}

\subsection{Formale Prinzipien des GitOps-Betriebsmodells}
\label{subsec:formale-prinzipien-gitops}

Das GitOps-Paradigma führt den deklarativen Infrastrukturansatz weiter und überträgt die Mechanismen der verteilten Versionsverwaltung auf den operativen Betrieb von Cloud-Native-Systemen. Um den Begriff wissenschaftlich und industriell eindeutig abzugrenzen, hat die \textit{OpenGitOps}-Arbeitsgruppe der \gl{cncf} vier verbindliche Kernprinzipien formuliert, die ein System erfüllen muss, um als GitOps-Architektur klassifiziert zu werden:\autocite{opengitops_principles_2021}

\begin{enumerate}
\item \textbf{Deklarativität (\textit{Declarative}):} Ein System, das über GitOps verwaltet wird, muss seinen gewünschten Zielzustand vollständig deklarativ ausdrücken.
\item \textbf{Versioniert und unveränderlich (\textit{Versioned and Immutable}):} Der gewünschte Zustand wird so gespeichert, dass Unveränderbarkeit, Versionierung und eine vollständige Versionshistorie erzwungen werden. In der Praxis übernimmt ein Git-Repository diese Funktion, wodurch jede Zustandsänderung durch einen kryptografisch signierbaren Commit-Hash (\gl{sha}) eindeutig identifizierbar und über ein \texttt{git revert} deterministisch umkehrbar ist.\autocite{beetz_gitops_2022,opengitops_principles_2021}
\item \textbf{Automatisch abgerufen (\textit{Pulled Automatically}):} Software-Agenten rufen die gewünschten Zustandserklärungen eigenständig und kontinuierlich aus der Datenquelle ab (\textit{Pull-Prinzip}). Im Gegensatz zum imperativen Push-Modell initiiert das Zielsystem die Verbindung nach außen, sodass keine Cluster-Zugangsdaten in externen Systemen exponiert werden müssen.\autocite{beetz_gitops_2022,opengitops_principles_2021}
\item \textbf{Kontinuierlich abgeglichen (\textit{Continuously Reconciled}):} Software-Agenten beobachten fortlaufend den tatsächlichen Systemzustand und unternehmen automatisierte Korrekturmaßnahmen, sobald eine Abweichung (\textit{Drift}) zum deklarierten Zustand festgestellt wird.\autocite{opengitops_principles_2021}
\end{enumerate}

Für die Untersuchung von Fehlerszenarien ist insbesondere das Zusammenspiel des dritten und vierten Prinzips von zentraler Bedeutung. Da der abgleichende Agent innerhalb des Clusters lokalisiert ist, die maßgebliche Wahrheitsquelle (das Git-Repository) jedoch außerhalb des Clusters liegt, entsteht eine permanente Abhängigkeit von der Netzwerkverbindung zwischen Control Plane und Versionsverwaltung.\autocite{beetz_gitops_2022,opengitops_principles_2021,ghosh_gitops_2025}

\section{Kontrolltheoretische Grundlagen: Das Operator-Pattern und die Reconciliation-Schleife}
\label{sec:operator-pattern-reconciliation}

Um das dynamische Verhalten von Kubernetes und darauf aufbauenden GitOps-Controllern wie ArgoCD unter Störeinflüssen präzise zu beschreiben, lässt sich deren Architektur auf Konzepte der Regelungstheorie (\textit{Control Theory}) zurückführen.\autocite{burns_borg_2016,lim_control_2009} In der Regelungstechnik überwacht ein Regler (\textit{Controller}) kontinuierlich die Regelgröße eines Systems, vergleicht diese mit einer vorgegebenen Führungsgröße (\textit{Setpoint}) und wirkt über eine Stellgröße korrigierend auf die Regelstrecke ein, um den Einfluss externer Störgrößen (\textit{Disturbances}) zu kompensieren.\autocite{beyer_sre_2016,lim_control_2009}

Kubernetes überträgt dieses Prinzip in Form der \textit{Reconciliation-Schleife} (\textit{Reconcile Loop}) auf die Softwaretechnik. Jeder Controller innerhalb der Kubernetes Control Plane arbeitet als geschlossener Regelkreis, der nach dem Muster \textit{Observe -- Diff -- Act} (Beobachten -- Vergleichen -- Handeln) operiert:\autocite{burns_borg_2016,dobies_operators_2020}

\begin{equation}
\Delta S(t) = S_{\text{desired}}(t) - S_{\text{actual}}(t)
\label{eq:reconcile-diff}
\end{equation}

Solange die Zustandsdifferenz $\Delta S(t) = \emptyset$ beträgt, befindet sich das System im eingeschwungenen Gleichgewichtszustand (\textit{Steady State} bzw. \texttt{Synced} und \texttt{Healthy}). Tritt eine Abweichung $\Delta S(t) \neq \emptyset$ auf -- sei es durch eine Aktualisierung des Soll-Zustands $S_{\text{desired}}(t)$ im Git-Repository oder durch einen Ausfall bzw. eine manuelle Veränderung im Cluster $S_{\text{actual}}(t)$ --, berechnet der Controller die erforderlichen \gl{api}-Operationen, um $\Delta S(t)$ wieder gegen Null zu führen.\autocite{burns_borg_2016,dobies_operators_2020}

Ein wesentliches Architekturmerkmal dieser Regelungsschleife ist die \textit{pegelgetriebene Triggerung} (\textit{Level-Triggered Reconciliation}) im Kontrast zu einer rein \textit{flankengetriebenen Verarbeitung} (\textit{Edge-Triggered}).\autocite{burns_borg_2016,dobies_operators_2020} Während ein flankengetriebenes System lediglich auf den unmittelbaren Zeitpunkt eines Ereignisses (etwa einen Webhook bei einem \texttt{git push}) reagiert, bewertet ein pegelgetriebenes System periodisch den gesamten aktuellen Zustand im Vergleich zum Zielzustand. Geht ein einzelnes Benachrichtigungsereignis aufgrund einer transienten Netzwerkstörung verloren, stellt die pegelgetriebene Architektur sicher, dass das System bei der nächsten erfolgreichen Zustandsabfrage dennoch zuverlässig konvergiert.\autocite{burns_borg_2016,dobies_operators_2020}

Das \textit{Kubernetes-Operator-Pattern} erweitert diesen nativen Regelungsmechanismus, indem domänenspezifisches Betriebswissen in eigene Controller gegossen und über \textit{Custom Resource Definitions} (\gl{crd}) direkt in die Kubernetes-\gl{api} integriert wird.\autocite{dobies_operators_2020} ArgoCD nutzt dieses Muster, indem es die Custom Resource \texttt{Application} definiert. Diese Ressource verknüpft eine externe Quelle (Repository-URL, Ziel-Revision und Verzeichnispfad) mit einem Ziel-Cluster und spezifischen Synchronisationsrichtlinien (wie \texttt{selfHeal} zur automatischen Korrektur von Cluster-Drift und \texttt{prune} zum Entfernen veralteter Ressourcen).\autocite{argocd_architecture_2026,dobies_operators_2020}

\section{Fehlertoleranz, Resilienz und Chaos Engineering in verteilten Systemen}
\label{sec:fehlertoleranz-resilienz-chaos}

Verteilte Systeme, die über Netzwerkgrenzen hinweg kommunizieren, unterliegen zwangsläufig partiellen Ausfällen. Um die Zuverlässigkeit solcher Architekturen wissenschaftlich zu bewerten, bedarf es einer klaren Abgrenzung der Begriffe aus der Verlässlichkeitsforschung (\textit{Dependability}) sowie etablierter Metriken zur Quantifizierung des Wiederherstellungsverhaltens.

\subsection{Fehlerkette (Fault -- Error -- Failure) und Resilienzbegriff}
\label{subsec:fehlerkette-resilienz}

Nach der grundlegenden Taxonomie von Avizienis et\,al. wird die Entstehung von Systemausfällen als dreistufige Kausalkette modelliert:\autocite{avizienis_dependable_2004} Ein \textit{Fault} (Fehlerursache bzw. Störung) bezeichnet die zugrundeliegende oder externe Ursache eines Fehlverhaltens, im Kontext dieser Arbeit etwa den partiellen Verlust von Netzwerkpaketen auf der Transportschicht. Führt dieser \textit{Fault} innerhalb des Systems zu einem abweichenden internen Teilzustand -- beispielsweise einem abgebrochenen Git-Abruf oder dem Übergang der ArgoCD-Zustandsmaschine in die Bedingung \texttt{ComparisonError} --, liegt ein \textit{Error} (Fehlzustand) vor. Erst wenn dieser interne Fehlzustand die nach außen sichtbare Dienstleistung des Gesamtsystems beeinträchtigt oder vollständig unterbricht, spricht man von einem \textit{Failure} (Ausfall).\autocite{avizienis_dependable_2004}

Der Begriff der \textit{Resilienz} (\textit{Resilience}) beschreibt vor diesem Hintergrund die Fähigkeit eines Systems, trotz des Auftretens von \textit{Faults} und resultierenden \textit{Errors} ein akzeptables Dienstniveau aufrechtzuerhalten (\textit{Graceful Degradation}) und nach dem Abklingen der Störung eigenständig sowie innerhalb einer definierten Zeitspanne wieder in den regulären Betriebszustand zurückzukehren.\autocite{beyer_sre_2016,avizienis_dependable_2004} In einem GitOps-System manifestiert sich Resilienz auf zwei Ebenen: Erstens darf ein Kommunikationsverlust zum externen Git-Repository die bereits im Cluster laufenden Produktiv-Workloads (\textit{Data Plane}) nicht beeinträchtigen (\textit{statische Resilienz}). Zweitens muss die \textit{Control Plane} nach Ende der Netzwerkstörung ohne manuellen Eingriff und ohne Ressourcenerschöpfung (wie \gl{oom}-Abstürze) wieder in den Zustand \texttt{Synced} und \texttt{Healthy} konvergieren (\textit{dynamische Resilienz}).\autocite{fogli_k3s_2021,sun_sieve_2022,argocd_architecture_2026}

\subsection{Zeitliche Resilienzmetriken: Verweildauer und Mean Time to Recovery (MTTR)}
\label{subsec:zeitliche-resilienzmetriken}

Zur quantitativen Erfassung der dynamischen Resilienz werden in der Systemevaluierung und im \textit{Site Reliability Engineering} (\gl{sre}) zeitbasierte Kennzahlen herangezogen.\autocite{beyer_sre_2016,nygard_release_2018} Im Rahmen einer kontrollierten Fehlerinjektion, die zum Zeitpunkt $t_{\text{fault\_start}}$ beginnt und zum Zeitpunkt $t_{\text{fault\_end}}$ endet, lassen sich zwei präzise abgegrenzte Zeitintervalle definieren:

\begin{enumerate}
\item \textbf{Transiente Verweildauer (\textit{Dwell Time} $T_{\text{dwell}}$):} Bezeichnet die gesamte Zeitspanne zwischen dem erstmaligen Eintritt des Systems in einen degradierten oder fehlerbehafteten Zustand ($t_{\text{degraded}}$) und der vollständigen Wiederherstellung des regulären Systemzustands ($t_{\text{recovered}}$):
\begin{equation}
T_{\text{dwell}} = t_{\text{recovered}} - t_{\text{degraded}}
\label{eq:dwell-time}
\end{equation}
\item \textbf{Mean Time to Recovery ($\text{\gl{mttr}}_{\text{post-fault}}$):} Im engeren messtechnischen Sinn eines kontrollierten Störungsexperiments beschreibt die Erholungszeit nach Störungsende die Differenz zwischen dem Zeitpunkt, an dem die externe Netzwerkstörung behoben ist ($t_{\text{fault\_end}}$), und dem Zeitpunkt, an dem die Control Plane den fehlerfreien Zielzustand (\texttt{Synced} und \texttt{Healthy} ohne Fehler-Conditions) wieder vollständig erreicht hat:\autocite{beyer_sre_2016,nygard_release_2018}
\begin{equation}
\text{\gl{mttr}}_{\text{post-fault}} = t_{\text{recovered}} - t_{\text{fault\_end}}
\label{eq:mttr}
\end{equation}
\end{enumerate}

Diese Differenzierung ist für die Bewertung von GitOps-Controllern essenziell: Während $T_{\text{dwell}}$ auch durch die Dauer der von außen aufgeprägten Netzwerkstörung bestimmt wird, isoliert $\text{\gl{mttr}}_{\text{post-fault}}$ die intrinsische Erholungsgeschwindigkeit der Control Plane, welche maßgeblich durch Polling-Intervalle, Timeout-Schwellenwerte und exponentielle Backoff-Algorithmen determiniert wird.\autocite{beyer_sre_2016,argocd_architecture_2026}

\subsection{Methodik des Chaos Engineerings}
\label{subsec:methodik-chaos-engineering}

Um die genannten Resilienzmetriken nicht erst bei unvorhersehbaren Produktionsausfällen passiv zu beobachten, sondern systematisch im Experiment zu bestimmen, wurde die Disziplin des \textit{Chaos Engineerings} entwickelt. Nach Basiri et\,al. bezeichnet Chaos Engineering das Durchführen gezielter, kontrollierter Experimente an einem verteilten System, um Vertrauen in dessen Fähigkeit aufzubauen, turbulenten und fehlerhaften Bedingungen im Betrieb standzuhalten.\autocite{basiri_chaos_2016}

Ein wissenschaftlich fundiertes Chaos-Experiment folgt dabei vier streng strukturierten Phasen:\autocite{basiri_chaos_2016} Zunächst wird der \textit{Steady State} (Normalzustand) des Systems anhand messbarer Ausgabemetriken (hier: Reconcile-Latenz, Synchronisationsstatus und Ressourcenverbrauch) als Referenzbasis (\textit{Baseline}) definiert. Darauf aufbauend wird die Hypothese formuliert, dass dieser Steady State auch bei Eintritt einer spezifischen Störung erhalten bleibt beziehungsweise nach deren Ende innerhalb definierter Grenzen wiederhergestellt wird. Anschließend werden reale Fehlervariablen -- wie Netzwerk-Paketverluste, Latenzen oder Serverabstürze -- kontrolliert in das System injiziert, während der \textit{Blast Radius} (Wirkungsbereich der Störung) strikt begrenzt und überwacht wird. Abschließend wird die Abweichung zwischen dem Verhalten unter Störung und der Baseline statistisch ausgewertet, um Schwachstellen in der Konvergenzlogik oder der Ressourcendimensionierung aufzudecken.\autocite{basiri_chaos_2016}

\section{Netzwerkschicht-Interaktionen und Betriebssystem-Einflussfaktoren}
\label{sec:netzwerkschicht-betriebssystem}

Da die Kommunikation zwischen einem GitOps-Controller und einem externen Git-Repository (z.\,B. GitHub) über \gl{https} (\gl{tcp}-Port 443) abgewickelt wird, wirkt sich ein Paketverlust auf der Netzwerkschicht unmittelbar auf das Zeitverhalten der höheren Protokollschichten aus. Gemäß der \gl{tcp}-Spezifikation (\gl{rfc} 9293) sichert das \textit{Transmission Control Protocol} eine zuverlässige Datenübertragung durch Sequenznummern, Empfangsbestätigungen (\gl{ack}) und zeitgesteuerte Neuübertragungen (\textit{Retransmissions}) ab.\autocite{eddy_rfc9293_2022}

Gehen Netzwerkpakete -- etwa durch Überlastung auf \gl{wan}-Strecken oder durch eine gezielte Injektion mittels des Linux-Kernel-Subsystems \texttt{netfilter}/\texttt{iptables} -- mit einer hohen Wahrscheinlichkeit $p$ verloren, greifen im \gl{tcp}-Stack des Betriebssystems gestaffelte Erholungsmechanismen:\autocite{eddy_rfc9293_2022,netfilter_iptables_2026} Beim \textit{Three-Way-Handshake} zum Verbindungsaufbau (\texttt{SYN} $\rightarrow$ \texttt{SYN-ACK} $\rightarrow$ \texttt{ACK}) sowie bei vollständigem Ausbleiben von Bestätigungspaketen tritt ein \textit{Retransmission Timeout} (\gl{rto}) ein. Nach jedem erfolglosen Übertragungsversuch verdoppelt der Linux-Kernel das \gl{rto}-Intervall gemäß einem exponentiellen Backoff-Verfahren.\autocite{eddy_rfc9293_2022} Bei einer hohen Paketverlustrate von beispielsweise $p = 0{,}80$ (80\,\% Paketverlust) führt bereits die sequenzielle Verkettung aus \gl{dns}-Auflösung, \gl{tcp}-Handshake, \gl{tls}-1.3-Schlüsselaushandlung und dem eigentlichen Git-Protokolltransfer (\texttt{git ls-remote} bzw. \texttt{git fetch}) dazu, dass einzelne Abfragen über viele Sekunden blockieren oder das applikationsseitige Timeout überschreiten.\autocite{argocd_architecture_2026,eddy_rfc9293_2022}

Auf Seiten der Control Plane hat dieses Verhalten direkte Auswirkungen auf den Ressourcenbedarf. Während blockierte Netzwerk-Goroutinen im Controller und im Repository-Server auf den Ablauf von \gl{tcp}- und \gl{http}-Timeouts warten, belegen die zugehörigen Datenstrukturen, Puffer und \gl{tls}-Kontexte weiterhin Arbeitsspeicher (\textit{Working Set Memory}). Auf kostenoptimierten Cloud- oder Edge-Knoten mit knapp bemessenem physischem \gl{ram} birgt dies ein methodisches Risiko für die Leistungsbewertung: Überschreitet der Speicherbedarf die Grenze des physischen Hauptspeichers, beginnt der Linux-Kernel, Speicherseiten auf den Sekundärspeicher auszulagern (\textit{Swapping} bzw. \textit{Paging}), oder beendet Prozesse über den \gl{oom}-Killer.\autocite{gregg_systems_2020,kernel_vmstat_2026} Da Zugriffe auf ein \gl{ssd}-basiertes Swapfile um Größenordnungen höhere Latenzen aufweisen als Zugriffe auf physischen \gl{dram}, würde ein unbemerktes Swapping die gemessenen Reconcile-Latenzen verfälschen (\textit{Messartefakt}). Eine wissenschaftlich valide Latenzmessung in ressourcenbeschränkten Umgebungen setzt daher zwingend den Nachweis voraus, dass während des gesamten Experiments keinerlei Speicherseiten ein- oder ausgelagert wurden ($\Delta \text{pswpin} = 0$ und $\Delta \text{pswpout} = 0$ in \texttt{/proc/vmstat}).\autocite{gregg_systems_2020,kernel_vmstat_2026}
